\chapter{Compact linear mappings and perturbations}
\setcounter{exercisepar}{0}

In this chapter, all vector spaces considered are vector spaces over a field K equal to R or C. The references to EVT, II made in this book concern in general the case K = R; for the case K = C, cf. EVT, II, p. 64 to 68.

A semi-normed space is a vector space E equipped with a semi-norm p and the topology defined by p (cf. EVT, II, p. 2). The unit ball of E, or of p, is the set of elements x of E such that \(p(x) \leqslant 1\) .

For every vector space E, we denote by \(1_{E}\) the identity mapping of E. If E, F and G are vector spaces and u: \(E \rightarrow F\) , \(v: F \rightarrow G\) are linear mappings, we sometimes denote by vu the linear mapping \(v \circ u\) from E into G.

Let E and F be topological vector spaces. A continuous linear mapping from E into F will also be called a morphism, and an endomorphism when F = E. A bijective linear mapping from E into F which is continuous and whose inverse is also continuous will be called an isomorphism and, when F = E, an automorphism. When E and F are semi-normed spaces, the notions of endomorphism, isomorphism and automorphism refer to the underlying topological vector space structure.

Given topological vector spaces E and F, we denote by \(\mathcal{L}^{\mathrm{f}}(\mathrm{E};\mathrm{F})\) the vector space of continuous linear mappings of finite rank from E into F. We also denote by \(\mathcal{L}^{\mathrm{f}}(\mathrm{E})\) the vector space \(\mathcal{L}^{\mathrm{f}}(\mathrm{E};\mathrm{E})\) .

We recall that a continuous linear mapping u from E into F is strict if it induces, by passage to the quotient, an isomorphism of E/ Ker(u) onto u(E) (TG, III, p. 16, def. 1); this is equivalent to saying that the image of every neighbourhood of 0 in E is a neighbourhood of 0 in u(E) (TG, III, p. 16, prop. 24).

